\documentclass[../main]{subfiles}
\begin{document}
\chapter{Vaporización y Gases reales.}

\img{images/11/vaporizacion.jpg}{Agua evaporándose a $100^\circ$ C}

\section{Ecuación de Van der Waals}

\subsection{Gases reales}

\img{images/11/fuerzas.jpg}{Interacción entre moléculas}

La ecuación de Van der Waals es una corrección de la ecuación de
gases ideales para describir el comportamiento de gases reales,
considerando las interacciones intermoleculares y el volumen finito
de las moléculas. Se expresa de la siguiente manera:

\[
  \left( P + \frac{a}{V_m^2} \right) \left( V_m - b \right) = RT
\]

Donde:
\begin{itemize}
  \item \( P \) es la presión del gas,
  \item \( V_m \) es el volumen molar,
  \item \( T \) es la temperatura,
  \item \( R \) es la constante universal de los gases,
  \item \( a \) y \( b \) son los coeficientes de Van der Waals.
\end{itemize}

Los coeficientes \(a\) y \(b\) tienen los siguientes significados:

\begin{itemize}
  \item \textbf{Coeficiente \( a \)}: Corrige las fuerzas de
    atracción ínter moleculares. A mayor \(a\), mayor es la atracción
    entre las moléculas.
  \item \textbf{Coeficiente \( b \)}: Representa el volumen ocupado
    por las moléculas del gas. A mayor \(b\), menor es el volumen
    disponible para el movimiento del gas.
\end{itemize}

Un \textit{gas real} o \textit{vapor real} es aquel en el que las
moléculas no solo ocupan espacio, sino que también interactúan entre
sí. La ecuación de Van der Waals ajusta el comportamiento ideal,
siendo especialmente útil en situaciones de bajas temperaturas o
altas presiones, donde las desviaciones del comportamiento ideal son
más notorias.

\actividad{Responder}
\begin{enumerate}
  \item Definir en que consiste el vapor húmedo.
  \item Definir el vapor seco.
  \item Explicar que es la saturación térmica.
  \item En que consiste la sobresaturación.
  \item ¿Qué es vapor saturado y seco?
  \item Realizar diagrama temperatura-entropía para el vapor saturado
    líquido-húmedo- saturado y seco. (campana).
  \item Explicar el funcionamiento de una máquina simple de vapor.
  \item Buscar tablas de equilibrio de vapor de agua a diferentes
    presiones y temperatura. Adjuntar al trabajo.
  \item ¿Qué es el título de un vapor y como se calcula? Explicar.
  \item Explicar la ecuación de estado de gases reales, coeficientes
    \emph{a y b}.
\end{enumerate}
\actividad{Resolver}
\begin{enumerate}
  \item Determinar la presión a que está sometido 1kg de amoníaco
    que, a la temperatura de 7°C, ocupa un volumen de 0,800m³.
  \item Determinar la presión a que están sometidos 10 moles de
    hidrógeno que a la temperatura de 127°C ocupan un volumen de 0,8m³.
  \item Calcular el peso de una masa de metano que posee un volumen
    de 0,5m³ a la presión de \SI{120}{\atm} y a la temperatura de 77C.
  \item Calcular la cantidad de calor necesaria para convertir 10 kg
    de agua a 15°C en "vapor húmedo" de titulo 0,8 a la presión de 5 atmósferas.
  \item Calcular la cantidad de calor necesaria para transformar 1kg
    de agua a 10°C, en vapor recalentado, a la presión de 8
    atmósferas y a la temperatura de 200°C, suponiendo que el calor
    específico medio del vapor recalentado valga $0,6\frac{kcal}{kg °C}$.
  \item Calcular la cantidad de calor necesaria para convertir 10 kg
    de agua a 20°C en "vapor húmedo" de titulo 0,9 a la presión de 6 atmósferas.
  \item Una masa de agua a la temperatura de 30°C y que pesa 20Kg es
    vaporizada a la presión constante de 16 atmósferas. Calcular la
    temperatura que alcanzará el vapor recalentado si durante todo el
    proceso se necesitaron 15000 kilo calorías y el calor específico
    medio del vapor es de 0,58 $\frac{kcal}{kg °C}$
  \item Determinar la energía interna de 10kg de agua a la
    temperatura de saturación de 179°C.
  \item Calcular la entalpía correspondiente a una masa de
    \emph{vapor de agua saturado y seco} que pesa 20kg y se encuentra
    a la temperatura de saturación 158,1C.
  \item Calcular la entalpía de 1kg de \emph{vapor de agua saturado
    húmedo} con título 0,85 cuya presión es de 14 atmósferas.
  \item Calcular la variación d energía interna necesaria para
    producir 1kg de \emph{vapor d agua saturado y húmedo} con titulo
    0,4 a la presión de 40 atmósferas.
  \item Calcular la entalpía de 1kg de \emph{vapor de agua saturado
    húmedo} con título 0,9 cuya presión es de 20 atmósferas.
  \item Calcular el peso de la parte de agua y de la parte de vapor
    seco correspondiente a 1m³ de vapor saturado húmedo de titulo
    0,65 que se encuentra a la temperatura de 151,1°C.
  \item Calcular el volumen, la entalpía y la energía interna de 1
    tonelada de vapor de agua, saturado húmedo con título 0,94 y a la
    presión de 18 atmósferas.
\end{enumerate}
\end{document}

